\section{The Philosophy of Chaos}

\begin{quotex}
Men love those enticements. The alternative is to purge one's mind of all vain opinions, desires, fears, anxieties, and so on. Men fear to do that, since they believe there will be nothing left, they will no longer have a self or an identity. On the contrary, in that clearing the Spirit will generate a new Self, the Logos, the opposite of the satanic chaos of the mind. \flright{\textit{Day of the Dead and Signs of Life}\footnote{\url{https://www.gornahoor.net/?p=5164}}}

Nothing under the sun is new, neither is any man able to say: Behold this is new: for it hath already gone before in the ages that were before us. \flright{\textsc{Ecclesiastes 1:10}}

And from the days of John the Baptist until now, the kingdom of heaven suffereth violence, and the violent bear it away. \flright{\textsc{Matthew 11:12}}

\end{quotex}
Being a review of \textit{Political Platonism}\footnote{\url{https://arktos.com/product/political-platonism/}} by Alexander Dugin.

\paragraph{Chaos}
This book is a collection of outtakes. To ruin a mystery, we turn to the end first, where Dugin explains:

\begin{quotex}
Our task is the construction of a philosophy of Chaos. 

\end{quotex}
Dugin has neglected to construct such a philosophy. Chaos is matriarchal, inclusive, anti-hierarchical, anti-violence, and so on; certainly, all ideas opposed by Logos. Unfortunately, Dugin gets off to a bad start when he asserts:

\begin{quotex}
Logos is based on the exclusion of Chaos, on the affirmation of a strict alternative to it. 

\end{quotex}
Logos does not seek the exclusion of Chaos. Rather it seeks mastery over Chaos. That is the masculine way that Dugin rejects. After rejecting some scientific misunderstandings of Chaos, Chaos begins to sound like the prime matter of Aristotle and the Scholastics. As such, it is an object of Thought, not of Being. You can show me a squirrel but you cannot show me a piece of Chaos: it has no Being.

He then claims that Chaos contains all possibilities. We addressed this same issue several years ago in \textit{Chaos in Mathematics}\footnote{\url{https://www.gornahoor.net/?p=4970}}. Chatin's Constant does indeed contain all possibilities, but only by accident. Moreover, only the Logos can recognize those pseudo-regularities in the digits of the Constant, although they are actually random.

Dugin gives us two choices: start at the beginning of philosophy or start at the end ``when the disappearance and liquidation of philosophy is occurring". Although Chaos is ``eternal", there is no beginning, i.e., we can't anticipate the beginning of the next cycle. Hence, we shall start at the End, since we can look back over the history of philosophy from the Greeks, through the Medievals, and into modernity.

We have documented the ``liquidation of philosophy" many times, e.g., from Hobbes, Bacon, Descartes, and so on. Anyone can see that. However, we encounter a serious problem whenever we try to penetrate to the prior period. The Medieval period is opaque to the modern mind, which cannot penetrate into a proper understanding of it. The texts can be analyzed, but the mindset is beyond mysterious. Only a remnant, born with a still archaic mind, can understand that period.

Hence, the fundamental premise that the philosophy of Logos contains the seeds of its own dissolution. Modernity is based on inferior rational thinking, whereas the Tradition of the Middle Ages is based on Being, an intuition of the whole, which is itself beyond logic, yet not illogical.

\paragraph{Tradition}
Dugin then tries to enlist \textbf{Rene Guenon} in support of his unfinished, or rather unstarted, project. This is a tactical error since there is no support from Guenon. Guenon accepts Aristotle and the Scholastics as representatives of Traditional metaphysics. Guenon rejects philosophy. On the other hand, metaphysics is eternally true and does not contain any seed of its dissolution. Moreover, once the metaphysical principles are realized, there is no way to return to error.

As Guenon explains, the unactualized possibilities are included in non-Being. The Void, which seems to be what Dugin means by Chaos, has no Being. Rather, the possibilities reside in the Mind of God as ideas or essences. To develop a philosophy that includes possibilities that can never manifest sounds like a fool's errand.

Nevertheless, Guenon agrees with Dugin about the decline of philosophy in modernity. Moreover, he also agrees that philosophy declines over time. However, that is not a logical decline. Specifically, it is not an object of Thought but rather a change of Being. There is no inner logic in true metaphysics that leads to Chaos, never mind a philosophy of Chaos.

If there is no inner logic within Tradition that leads to Chaos, then there must be transcendent forces that bring such changes about. This is documented in \textit{Angels and Demons}\footnote{\url{https://www.gornahoor.net/?p=11086}}. So, if European philosophy has failed, that is not a result of bad logic. Rather, it is a symptom of the effects of diabolical forces infecting the European mind.

\paragraph{Kali Yuga}
Chaos, as described by Dugin, has much in common with the time of the Kali Yuga. Specifically, it is undifferentiated. It rejects hierarchy. There is no need to repeat what has previously been described.

\paragraph{Further Openings}
There are some deeper ways to deal with the idea of Chaos.

The Cabala describes certain world periods as Chaos. In particular, the three world ages—in Christian terms, the ages of the Father, the Son, and the Holy Spirit—correspond to Chaos, Logos, and the Messiah. These refer to Being, not abstract Thought.

Jacob Boehme, Meister Eckart, inter alia, have had intuitions of non-being beyond Being.

The writings of Vladimir Solovyov on the Divine Sophia are helpful. Sophia, as the Soul of the World, keeps Chaos in check insofar as she is immaculate. That gives rise to the Logos.

The point is that there is no philosophical solution to the problem of Chaos. Only a metaphysical one.

This is the second attempt to deal with the philosophy of chaos. A third is probably necessary, because something seems to be missing.



\flrightit{Posted on 2019-11-03 by Cologero }

\begin{center}* * *\end{center}

\begin{footnotesize}\begin{sffamily}



\texttt{Jack on 2019-11-04 at 08:04 said: }

Pro-chaos, anti-Logos. Alright then, good to know I need not waste any time reading anything by Mr. Dugin. Plenty of other good books waiting, nice to check some off the list.


\hfill

\texttt{JULIANO P CORREA on 2019-11-04 at 09:02 said: }

Your points are some of the conclusions I had that made me return to catholicism, and later on reject esotericism as a whole. Dugin's idea come from his Heideggarian (Heiddeger thinks the problem of philosophy is in the very thinking of the Logos) thought that he won't give up, and the result is he become a adherent of the path of the left hand. Your differences with him come down to that, he belives the left hand only is correct, and you that right hand is correct. It seems that Guénon do not like the left hand a lot, but he still identify scholastics as inferior to what he call complete metaphysics, since for him when the absolute God is indetified with Being, like catholic, that is inferior system compared to when it goes beoynd it, like advaita vedanta. This is clear in his book on the vedanta and other texts. The difference comes down to if one thinks that non-being ``exists" depending on being as its privation only, or if being comes from non-being. The first is what classical religions say, the second is what esotericism says. The problem with right hand esotericism is simply that it automatically leads to left hand esotericism as it's conclusion.

One discovers than that actually all esotericism is flawed, and everything the church condemned as heresy really is heresy, including Gioacchino da Fiore, Sovyov, and the other ``christian gnostics".


\hfill

\texttt{Cologero on 2019-11-08 at 09:04 said: }

Mr. Correa:

This kind of broad critique, with little idea density, is what we oppose. Instead, we have provided specific examples from multiple sources, over and over again. Some may be misleading, others correct. If you want to make your ideas clear, or engage in some sort of dialog, then choose some \emph{specific} point made in this post and offer your objections or suggestions.

BTW, this is what Guenon wrote of Aristotle in \emph{Introduction to Hindu Doctrine}. Hardly a knee pain, except possibly for you.

\begin{quotex}
They do not in fact appear to have introduced any Hindu Ainfluences, except that contained in the logic of Aristotle—to which we have already alluded in connection with his syllogism—and also in the metaphysical part of the same philosopher's work, in which it is possible to point to intellectual affinities with India far too close to be purely accidental. 

\end{quotex}

\hfill

\texttt{Hafez Feizi on 2019-11-10 at 14:00 said: }

Cologero, I believe you are showing signs of intemperance for shooing people away the moment they voice some inkling of disagreement. After all the man you just banned is a fellow Catholic Christian, not a neo-pagan or atheist with whom discussion is not possible since they do not share the common grounds for understanding. Furthermore, Juliano was not exhibiting an attitude of toxic arrogance or dismissiveness that you commonly see under the tyranny of opinion prevailing on internet comment sections. 

To be fair I can sense that you find it annoying to read a tirade loosely directed against everything laid out on this site without focusing on the particular article in question in a very generalized and broad headed manner. However I do believe his critique, from a Christian standpoint, is not without merit. You should at least give him a chance to defend himself and limit his disagreements to the topics at hand in the future.

Now without agreeing nor disagreeing with the points that he raised, I will state some observations of my own: 

What do we mean by esoteric? To my understanding, it is something that comes from beyond the human ability to directly sense (e.g. objects) or intellectually apprehend (energy). Sensible, phenomenal objects are not esoteric. God is not some sort of alien or energy lifeform like today's popular science loves to picture him. This leaves us with the numinal realm of thought and idea forms. To circumscribe the divine in the metaphysics of Being is to restrict it in an entirely impersonal framework of understanding. One of my problems with traditional metaphysicians is that nothing they write explains to me the intensely personal way in which Trinity manifests itself and communicates with the human world through history as described in the Biblical Scriptures. In some sense metaphysics ``restricts" God to behaving in an intellectually coherent and to some extent, predictable way characteristic of a systemic branch of reasoning. There is no shortage of ``esoteric" doctrines which treat the metaphysical realm as nothing but a collection of impersonal laws, numinal powers and efficient levers that can be rationally assessed and used to influence the world around, The question then boils down to how we can gain power over these forces to willfully effectuate their purpose on the natural world (not excluding our internal realm, in which the occult arts are primarily directed anyway), thus becoming ``like gods" ourselves. The purpose of esoteric discipline then becomes a set of rationally accessible alchemic guidelines and rites to transform man into something other than what he is created to be, thereby replacing the Creator's will for mankind's destiny with that of our own, personal making. It was precisely this type of doctrine that Juliano was perfectly justified in denouncing. 

Therefore if we are to accept a Catholic, Christian understanding of God, then He has to come from something above and beyond this. 

That is why prayer, regular confession, eucharist, asceticism and monasticism are indispensable to the life of the Church. Without them, the Christian existence would have no meaning, no point of reference. Juliano again, was right in saying that morality, understood in its proper sense, is central to the personal mystical communion with God, not esoteric disciplines of any sort, which might be helpful to some degree in pointing out the right orientation, but are inevitably misleading in and of themselves. Our egotistic desires, not only simple ones such as greed or lust for power, but also hidden, subtle impulses to yearn for the mystical union (which in truth is the origin of all sin) that we do not even readily recognize as such, are the primary reason why God is not manifest in the world in such a way that everyone can openly recognize Him. This is why even the most spiritually enlightened of us have to live our lives in total spiritual darkness, not being able to see or sense His presence around us, staving off the prevailing atmosphere of disbelief and nihilism on grounds of pure faith in something that is essentially missing in this carnival of appearances and smoke. But in no way should ``missing" point to the realm of ``non-being" or something of the like. And neither should we seek after some kind of substitute under the purpose of following the breath of the spirit (which I can assure you wouldn't have in this case anything ``Holy" attached to it). 

Another problem with esotericism is the tendency to create differences and distinctions according to the rank in which one has advanced in the doctrine. Now don't get me wrong, in general I'm very pro-elitist. Unqualified equality is in my view, an evil thing. But all matters of hierarchy are related with the utilitarian organisation of social order, and as such are not, in my view, part of the role in the realm of the spirit as God intended for humans to occupy. There is an angelic hierarchy in heaven, but to my understanding, humans are not a part of it. Otherwise we would just be lousy versions of Angels, which I doubt Christ would incarnate as to bother saving. Furthermore God never creates something deliberately meant to be inferior, lower versions of already existing things, as that would imply humans to be nothing but. 

Now this kind of esoteric hierarchy can allow one to essentially say: ``I've practiced this doctrine for longer than you, I committed the rites and sacrifices with greater attention and intensity, I have had better success in spiritual matters, therefore I'm placed on an ontologically different level from you." This is precisely the kind of attitude that the Gospel clearly cuts through when describing the Pharisees. The latter obeyed the law to the letter, prayed for hours a day and performed all the ritual cleansing rites. Yet Christ preferred the company of vagrants, outcasts and prostitutes to them. I wonder why that is? 

In one of your posts you state that one of the purposes of life is to acquire esoteric knowledge. I'm not sure if this is a well thought out statement. It can never be the goal of any Christian to acquire anything, except God's grace, and I don't see how that requires any knowledge of the occult. What esoteric practices did Mary commit herself to?


\hfill

\texttt{Cologero on 2019-11-10 at 15:44 said: }

You are doing better than Juliano, but not much better. I presume you know what a ``straw man" is; this ``esoteric hierarchy" — where exactly is it? — is not something I have promoted. At least you pointed to a sentence:

\begin{quotex}
one of the purposes of life is to acquire esoteric knowledge 

\end{quotex}
I don't believe I wrote that, and I certainly don't believe it. So I would have to revisit that quote and modify it or put it in context. From my perspective you have advanced the conversation.

I lost all respect for Juliano when he backtracked on his assertion regarding Aristotle. A man would have acknowledge his error and moved on. Perhaps in this way:

If Aristotle was close to one of the Hindu schools, then certainly Thomas Aquinas would have been, even if by accident. Hence, it could be worthwhile to investigate that school with the intent of deepening our understanding of the issues raised by Thomas. Is that so far fetched? Consider this quote from \textbf{Jacques Maritain}, certainly no esoterist from his book \emph{Preface to Metaphysics}: 

\begin{quotex}
In view of all this we are shocked if we are told of a knowledge which applies to-day the same fundamental concepts, the same principles as in the days of Shankara, Aristotle or St. Thomas. 

\end{quotex}
So do you super-correct Catholics still want to object to my use of Shankara? I am just finishing the job that Maritain never even started.

As for your many misunderstandings, I have no use for them. If you are no egalitarian, then you should concede that you are the more ignorant partner in this conversation.


\hfill

\texttt{Cologero on 2019-11-11 at 11:27 said: }

This is from Death and the Real I: \url{https://www.gornahoor.net/?p=11320}

\begin{quotex}
Acquire esoteric knowledge

So this quote exists nowhere on this site: ``one of the purposes of life is to acquire esoteric knowledge"

That suggestion was only for those who want to do the death mediation exercise.

Yet another example of those who cannot even read a simple text, yet come here to complain about our Terms of Service.

\end{quotex}

\hfill

\texttt{Hafez Feizi on 2019-11-11 at 16:17 said: }

It's not in my interest to argue with you Cologero. Nothing beneficial will come from it. I concede that on a closer look, I was mistaken in presuming that you claimed ``acquiring esoteric knowledge" is the primary means towards attaining gnosis. I would sound very silly if I found anything wrong with such acquisition in and of itself, since it's the very reason why I find the resources provided in this website valuable. 

The substance of my previous post was not even chiefly aimed in your direction, but rather to give voice to my concerns about misguided spiritual development. Far too many times we have been mislead by esoteric doctrines and figures which claim to restore traditional principles but subtly belay an empirical materialistic mindset (frequently in form of so called ``spiritual materialism"). I do believe some of these men acted in good faith and were admirable figures (such as Baron Evola, Rudolf Steiner, Gurdjieff and others) but in general, their teachings could lead to apostasy and heresy if accepted without condition. It can be extremely difficult to discern truth from falsehood without the guidance of dogma and two millenia of passed down Apostolic tradition. This is why I also strongly contend that it is dangerous and misleading to undertake a serious study in matters of the spirit and the divine without the prior formation of a spiritual bond with said tradition (both on a mental and spiritual level). This would also throw a wrench in the mouths of those critics who accuse you of fraud for promoting ``Romish" Christianity while at the same time upholding esoteric spiritual practices that it supposedly never openly endorsed.


\hfill

\texttt{Cologero on 2019-11-11 at 16:42 said: }

The comment section is not intended to be your own soapbox. I suggest that you create your own blog in which you can share your concerns. I promise to link to it.

But for now, you are under moderation.


\hfill

\texttt{Matt on 2019-11-12 at 00:25 said: }

Hafez and Juliano's comments serve as a helpful reminder to those who come across this site. No one is compelling you to frequent Gornahoor. If you are ultimately bothered by the content offered, it is probably best that you move on. It is not required for one's salvation. But with that said, some lingering misconceptions need to be cleared up. 

This site's aim is not to slavishly repeat what was said by Guenon – or any of the other pegged ``traditionalists" for that matter. It is not a neo-scholastic manual that inputs an answer from Guenon or Evola on a given question. Gornahoor has engaged the works of these writers, noting their insights, but also pointing out where they missed the mark. More importantly, those men, and the other individuals described as traditionalists, are not the only people whose works this site has explored. There have been dozens of names that have appeared in bold font in the site's posts throughout its existence. Many of those names are the Church Fathers and the saintly mystics. The reason for that is to point one in the direction of their works so that you can read them for yourself and start the process of trying to attune oneself to the state of consciousness that they were in when they wrote those works.

The importance of morality and virtue (the power to do what is right) has never been neglected here. I've lost count how many times Lorenzo Scupoli's Spiritual Combat has been referenced in posts.

Finally, the importance – nay, the necessity – of grace has never been neglected. Pelagius' heresy is not getting a second look here. We are not taking our cue from Valentinus and the other gnostics. Clement of Alexandria said it best when he described gnosis as illuminated faith.


\hfill

\texttt{Matt on 2019-11-12 at 00:46 said: }

``But all matters of hierarchy are related with the utilitarian organisation of social order, and as such are not, in my view, part of the role in the realm of the spirit as God intended for humans to occupy."

The consensus of the Fathers and Doctors of the Church would beg to differ. If you're a Catholic, you should take heed to what they have written on that.

Deification is at the heart of the Church's understanding of salvation. To be made God-like is no mere flowery metaphor. There are ontological implication that follow from this doctrine. There are a number of posts on this site that tease out those implications. But you don't have to read them, since you can read the writings of the saints and mystics that also show those implications.


\end{sffamily}\end{footnotesize}
